

\section{Програмна реалізація аналітичного комплексу}

\subsection{Загальна структура програмного комплексу}

Програмна частина системи розроблена з дотриманням принципів модульної декомпозиції,
коли кожен незалежний алгоритмічний крок інкапсулюється в окремому Python-модулі з
чітко визначеним публічним інтерфейсом. Такий підхід дозволяє замінювати окремі
компоненти (наприклад, ваги нейромережевої моделі або параметри \gls{IMM}-фільтра)
без зачіпання решти системи, а також значно спрощує одиничне тестування математичних
моделей.

Кореневою директорією проекту є \texttt{diploma\_project/}, у якій виділяються чотири
функціональні підкаталоги:
\begin{itemize}
    \item \texttt{ball\_3d\_analysis/} --- алгоритмічне ядро системи: модулі трекінгу,
    фільтрації, тривимірного позиціонування та трансформації вимірювань;
    \item \texttt{training\_models/} --- сценарії тренування моделі \gls{YOLO}26,
    конфігурація аугментацій навчальної вибірки, ваги навчених моделей;
    \item \texttt{parsing\_datasets\_scripts/} --- допоміжні сценарії обробки набору
    даних: розбиття на тренувальну, валідаційну та тестову вибірки, конвертація
    форматів анотацій;
    \item \texttt{masking\_scripts/} --- утилітарні засоби препроцесингу: інтерактивний
    калібратор діапазонів HSV для виокремлення м'яча, сценарій розпарсювання сирих
    відеопотоків на окремі кадри з анотаціями.
\end{itemize}

% TODO: створити файл sections/pictures_schemes/chapter_3/3_1_module_dependency_graph.png
% та розкоментувати цей блок
%\begin{figure}[!htbp]
%    \centering
%    \label{fig:module_dependency_graph}
%    \includegraphics[width=0.85\textwidth]{sections/pictures_schemes/chapter_3/3_1_module_dependency_graph}
%    \caption{Граф залежностей між модулями програмного комплексу}
%\end{figure}

Граф залежностей між модулями реалізує строгу односпрямовану структуру без циклічних
імпортів: модуль \texttt{IMMTracker} залежить від \texttt{IMM\_UKF}, який, у свою чергу,
залежить від \texttt{get\_court\_position\_methods}. Такий каскадний дизайн дозволяє
тестувати кожен модуль ізольовано, починаючи з найнижчого рівня геометричного
позиціонування.

\textbf{Обґрунтування вибору мови програмування.} Як основна мова реалізації обрано
Python 3.10 з огляду на такі чинники:
\begin{itemize}
    \item Багата екосистема готових наукових бібліотек (NumPy, SciPy) реалізує
    необхідні матричні операції з продуктивністю, наближеною до Fortran, завдяки
    зв'язуванню з BLAS/LAPACK.
    \item Прямий API-доступ до фреймворку Ultralytics для тренування та інференсу
    моделей \gls{YOLO} зі стандартною агрегацією метрик.
    \item Висока зрілість бібліотеки OpenCV (cv2), що містить готову реалізацію
    \texttt{cv2.solvePnP} та \texttt{cv2.Rodrigues}, які становлять ядро геометричного
    модуля.
    \item Прозорість та читабельність математичних формул при перенесенні з наукових
    публікацій безпосередньо в код.
\end{itemize}

Альтернативний варіант реалізації обчислювально-критичних частин на C++ через
прив'язку pybind11 був розглянутий і відхилений, оскільки профілювання показало, що
розроблений конвеєр виконується зі швидкістю $46-52$ FPS на типовій конфігурації
обладнання (RTX 4070, 32 ГБ RAM), що повністю задовольняє вимоги реального часу при
вхідній частоті відеопотоку $30-50$ FPS.

\subsection{Модуль геометричного зворотного проектування}

Модуль \texttt{get\_court\_position\_methods.py} реалізує два публічні API-методи, що
формують підвалини тривимірного відновлення координат об'єкта зі статичного
монокулярного відеозйомного засобу.

\textbf{Функція \texttt{calibrate\_camera(pts\_3d, pts\_2d, camera\_matrix, dist)}.}
Виконує одноразову калібрувальну процедуру камери на основі попередньо заданих
реперних точок розмітки волейбольного майданчика. Вхідні параметри:
\texttt{pts\_3d} --- масив форми $N \times 3$ з фізичними координатами реперних
точок у системі координат корту; \texttt{pts\_2d} --- масив форми $N \times 2$ з
відповідними піксельними координатами цих точок у кадрі; \texttt{camera\_matrix} ---
матриця внутрішніх параметрів $K$ розмірності $3 \times 3$; \texttt{dist} --- вектор
коефіцієнтів радіальної дисторсії об'єктива (приймається нульовим при умові
використання якісного об'єктива без помітних спотворень).

Внутрішня логіка функції включає виклик \texttt{cv2.solvePnP} з прапором
\texttt{SOLVEPNP\_ITERATIVE}, що активує ітераційний метод Левенберга-Марквардта.
Отримана вектор-обертова представлення $\vec{r}$ перетворюється на матрицю $R \in
SO(3)$ за формулою Родрігеса (\texttt{cv2.Rodrigues}), після чого розраховується
фізична позиція оптичного центру камери в світі: $\vec{C} = -R^T \vec{t}$. Функція
повертає кортеж $(R, \vec{t}, \vec{C})$, що повністю описує екстринсичні параметри
оптичної системи.

\textbf{Функція \texttt{get\_3d\_position(u, v, bbox\_w, K, R\_matrix, camera\_pos,
ball\_diameter)}.} Виконує зворотне проектування центру обмежувальної рамки м'яча у
тривимірну систему координат залу. Алгоритм складається з трьох послідовних кроків:
\begin{enumerate}
    \item Обчислення глибини $Z_c$ за формулою $Z_c = (f_x \cdot D_{real}) / w$, де
    $w$ --- ширина рамки в пікселях, $D_{real}$ --- стандартний діаметр м'яча 0.21
    м, $f_x$ --- горизонтальна фокальна відстань.
    \item Нормалізація піксельних координат у вектор напрямку променя в системі
    координат камери: $\vec{d}_{cam} = [(u - c_x)/f_x,\, (v - c_y)/f_y,\, 1]^T$.
    \item Трансформація напрямкового вектора у глобальну систему координат корту
    та обчислення кінцевої точки: $\vec{P}_{world} = \vec{C} + Z_c \cdot R^T
    \vec{d}_{cam}$.
\end{enumerate}

Усі константи, специфічні для конкретного запису (фізичні координати кутів корту,
піксельні координати цих кутів у кадрі, матриця $K$), винесені на рівень модуля та
ініціалізуються при імпорті. Це дозволяє повторно використовувати функцію для
кожного нового кадру без накладних витрат на повторну калібровку.

\subsection{Модуль ядра фільтрації \texttt{IMM\_UKF}}

Модуль \texttt{IMM\_UKF.py} є найбільш великим і алгоритмічно насиченим компонентом
системи (близько 680 рядків коду). Він реалізує три ключові класи та шість функціональних
методів, що утворюють базис для побудови взаємодіючих множинних моделей з
сигма-точковим калмановим фільтром.

\textbf{Клас \texttt{UnscentedKalmanFilter}.} Реалізує канонічну версію
сигма-точкового калманового фільтра з масштабованим перетворенням Ван дер Мерве.
Конструктор приймає набір параметрів: \texttt{name} --- ім'я-ідентифікатор фільтра
(для діагностики); \texttt{dim\_x = 6} --- розмірність вектора стану;
\texttt{dim\_z = 3} --- розмірність вектора вимірювання; \texttt{dt} --- крок
дискретизації за часом; \texttt{fx} --- функція переходу стану (одна з
\texttt{fx\_ballistic}, \texttt{fx\_hit}, \texttt{fx\_bounce}); \texttt{hx} ---
функція вимірювання (за замовчуванням виокремлює координати $[x, y, z]$ з вектора
стану); \texttt{points} --- об'єкт класу \texttt{MerweScaledSigmaPoints}.

Публічними методами класу є:
\begin{itemize}
    \item \texttt{predict(dt=None, UT=None, fx=None)} --- виконує крок прогнозу,
    обчислюючи апріорну оцінку стану $\hat{x}_{k|k-1}$ та її коваріацію
    $P_{k|k-1}$. Внутрішньо викликає \texttt{compute\_process\_sigmas} для генерації
    нової сім'ї сигма-точок, прогнозує кожну з них через нелінійну функцію
    \texttt{fx}, після чого застосовує сигма-точкове перетворення для відновлення
    моментів першого та другого порядку.
    \item \texttt{update(z, R=None, UT=None, hx=None)} --- виконує крок корекції за
    наявним вимірюванням $z$. У випадку $z = \text{None}$ (відсутність детекції)
    фільтр зберігає апріорний стан незмінним та встановлює \texttt{likelihood = 1.0},
    що забезпечує стабільність роботи \gls{IMM} при пропусках. У штатному режимі
    обчислюється підсилення Калмана $K = P^{xz} S^{-1}$, після чого оновлюються
    стан та коваріація.
    \item \texttt{cross\_variance(x, z, sigmas\_f, sigmas\_h)} --- допоміжний
    метод для обчислення перехресної матриці коваріації стану та вимірювання.
    \item \texttt{compute\_process\_sigmas(dt, fx=None, **fx\_args)} --- генерує
    набір 13 сигма-точок навколо поточної оцінки стану та пропускає кожну з них
    через функцію переходу.
\end{itemize}

\textbf{Клас \texttt{IMMEstimator}.} Реалізує оцінювач взаємодіючих множинних
моделей, що паралельно підтримує три екземпляри \texttt{UnscentedKalmanFilter} і
комбінує їхні оцінки за принципом байєсівської змішки. Конструктор приймає
\texttt{filters} --- список з трьох фільтрів; \texttt{mu} --- вектор початкових
ймовірностей моделей $[\mu_1, \mu_2, \mu_3]$; \texttt{M} --- матрицю Маркова
переходу між режимами.

Ключові методи:
\begin{itemize}
    \item \texttt{predict(u=None)} --- виконує цикл змішування початкових станів за
    вагами $\omega_{ij}$, що враховують ймовірність переходу між режимами, після
    чого ініціює прогноз для кожного з трьох \gls{UKF}.
    \item \texttt{update(z)} --- проводить оновлення кожного фільтра за вимірюванням
    $z$, зчитує функції правдоподібності кожної моделі, перераховує ваги $\mu$ за
    формулою $\mu_i \leftarrow \bar{c}_i \cdot \Lambda_i / \sum_j \bar{c}_j \Lambda_j$
    і обчислює змішану оцінку стану та її коваріацію.
    \item \texttt{\_compute\_mixing\_probabilities()} --- внутрішній метод для
    обчислення матриці змішування $\omega_{ij} = (\pi_{ij} \cdot \mu_i) / \bar{c}_j$.
    \item \texttt{\_compute\_state\_estimate()} --- внутрішній метод для відновлення
    змішаного стану $\hat{x} = \sum_i \mu_i \hat{x}_i$ та коваріації
    $\hat{P} = \sum_i \mu_i (\hat{P}_i + (\hat{x}_i - \hat{x})(\hat{x}_i - \hat{x})^T)$.
\end{itemize}

\textbf{Клас \texttt{MerweScaledSigmaPoints}.} Реалізує генератор детермінованих
сигма-точок за схемою Ван дер Мерве з трьома параметрами тонкої настройки
($\alpha$ --- розкид точок навколо середнього, $\beta$ --- врахування виду розподілу,
$\kappa$ --- допоміжний масштабний коефіцієнт). Основний метод
\texttt{sigma\_points(x, P)} обчислює корінь Холецького з матриці $(n + \lambda) P$
та генерує $2n + 1 = 13$ точок навколо вектора стану.

\textbf{Функції моделей руху.} Модуль містить три варіанти нелінійної функції
переходу стану, кожна з яких відповідає одному з режимів \gls{IMM}:
\begin{itemize}
    \item \texttt{fx\_ballistic(x, dt)} --- балістична модель з квадратичним
    аеродинамічним опором та гравітацією, що діє вздовж вертикальної осі $Y$.
    Швидкості кліпуються до $|v| \leq 40$ м/с для запобігання вибуху коваріації
    при сильно зашумлених початкових даних.
    \item \texttt{fx\_hit(x, dt)} --- лінійна модель сталої швидкості, що
    застосовується для коротких інтервалів навколо моментів ударних взаємодій,
    коли прискорення є розривним і не може бути передбачене.
    \item \texttt{fx\_bounce(x, dt)} --- модель кінематичного відскоку від
    підлоги, що інвертує вертикальну швидкість з коефіцієнтом реституції
    $\varepsilon = 0.75$ та згасає горизонтальні компоненти з коефіцієнтом
    тертя $\mu = 0.85$.
\end{itemize}

\textbf{Допоміжні функції.} \texttt{unscented\_transform(sigmas, Wm, Wc, noise\_cov)}
виконує власне сигма-точкове перетворення, повертаючи зважене середнє та зважену
коваріацію набору точок. \texttt{Q\_discrete\_white\_noise(dim, dt, var)} генерує
дискретну матрицю білого шуму для моделі постійної швидкості, що використовується
як блок матриці шуму процесу $Q$. \texttt{get\_dynamic\_transition\_matrix(y, vy)}
повертає поточну матрицю переходу режимів \gls{IMM} з урахуванням вертикальної
координати об'єкта та знаку вертикальної швидкості, реалізуючи адаптивну поведінку,
описану у підрозділі 2.4. \texttt{measurement\_transform(prediction\_data, K,
R\_matrix, camera\_pos, ball\_diameter)} є обгорткою, що приймає словник з
результатами детекції \gls{YOLO} ($x\_pos$, $y\_pos$, $w\_box$) та повертає вже
готове вимірювання $z = [X, Y, Z]^T$ у системі координат корту.

\subsection{Модуль трекінгу \texttt{IMMTracker}}

Модуль \texttt{IMMTracker.py} реалізує два класи високого рівня, що інкапсулюють
повний життєвий цикл відстеження об'єкта в часі.

\textbf{Клас \texttt{Track}.} Представляє один окремий просторовий трек об'єкта.
Поля класу:
\begin{itemize}
    \item \texttt{track\_id} --- унікальний цілочисельний ідентифікатор треку, що
    зберігається протягом усього часу життя об'єкта;
    \item \texttt{imm} --- екземпляр \texttt{IMMEstimator}, що відповідає за
    оцінювання поточного стану;
    \item \texttt{state} --- стан життєвого циклу (одне зі значень \textit{Tentative},
    \textit{Confirmed}, \textit{Deleted});
    \item \texttt{time\_since\_update} --- лічильник кадрів від останнього успішного
    оновлення треку детекцією;
    \item \texttt{hits}, \texttt{hit\_streak} --- лічильники сумарної та послідовної
    кількості успішних оновлень;
    \item \texttt{h\_matrix} --- матриця спостереження розмірності $3 \times 6$, що
    проектує 6-вимірний стан на 3-вимірне просторове вимірювання.
\end{itemize}

Поведінкові методи:
\begin{itemize}
    \item \texttt{predict()} --- оновлює адаптивну матрицю Маркова через виклик
    \texttt{get\_dynamic\_transition\_matrix} і виконує прогноз \gls{IMM}-оцінювача
    на один крок дискретного часу.
    \item \texttt{update(z)} --- виконує оновлення стану новою тривимірною
    координатою. Скидає лічильник \texttt{time\_since\_update}, збільшує
    лічильники \texttt{hits} і \texttt{hit\_streak}. Якщо стан треку був
    \textit{Tentative} і \texttt{hits} досягло порогу $\geq 3$, переводить трек у
    стан \textit{Confirmed}.
    \item \texttt{mark\_missed()} --- викликається, коли трек не отримав асоційованої
    детекції у поточному кадрі. Обнуляє \texttt{hit\_streak} та запускає крок
    екстраполяції стану \gls{IMM}-оцінювача без вимірювання.
    \item \texttt{get\_mahalanobis\_distance(z, R\_matrix)} --- обчислює квадрат
    статистичної відстані Махаланобіса між проекцією поточного апріорного стану
    треку на простір вимірювань та новою детекцією $z$. Повертає велике значення
    ($10^5$) у випадку виродженої матриці інновації.
\end{itemize}

\textbf{Клас \texttt{IMMTracker}.} Реалізує оркестратор колекції активних треків на
сцені. Поля класу:
\begin{itemize}
    \item \texttt{dt} --- крок дискретизації за часом ($1/\text{FPS}$);
    \item \texttt{max\_age = 50} --- максимальна кількість кадрів без оновлення, після
    якої трек переводиться у стан \textit{Deleted};
    \item \texttt{min\_hits = 3} --- мінімальна кількість послідовних підтверджень
    для переходу у стан \textit{Confirmed};
    \item \texttt{gating\_threshold = 11.34} --- поріг стробування за квадратом
    відстані Махаланобіса (відповідає $\chi^2$-розподілу з 3 ступенями свободи при
    довірчому інтервалі 99\%);
    \item \texttt{tracks} --- список активних треків;
    \item \texttt{next\_id} --- лічильник для генерації нових ідентифікаторів треків.
\end{itemize}

Основним публічним методом є \texttt{update(detections\_3d)}, що приймає список
тривимірних детекцій поточного кадру та виконує весь цикл асоціації, оновлення та
керування життєвим циклом (детально описаний у підрозділі 2.5 у вигляді псевдокоду).
Метод \texttt{get\_confirmed\_tracks()} повертає лише ті треки, що знаходяться у
стані \textit{Confirmed}, що використовується для подальшої візуалізації та
накопичення статистики траєкторії.

\subsection{Модуль навчання моделі \gls{YOLO}26}

Сценарій \texttt{training\_models/model\_train\_remote.py} реалізує повний цикл
тренування моделі детекції на віддаленому сервері з GPU. Логіка сценарію розділена
на чотири логічні блоки:
\begin{enumerate}
    \item \textbf{Завантаження базової моделі.} Через API
    \texttt{ultralytics.YOLO("yolo26m.pt")} ініціалізується модель середньої
    розмірності з попередньо натренованими вагами на COCO.
    \item \textbf{Конфігурація аугментацій.} Створюється pipeline бібліотеки
    Albumentations з наступним набором перетворень: \texttt{MotionBlur(blur\_limit=7,
    p=0.4)} для моделювання розмиття від руху; \texttt{GaussNoise(var\_limit=(10,
    50), p=0.3)} для імітації шуму матриці; \texttt{CLAHE(clip\_limit=2.0, p=0.4)}
    для підсилення локального контрасту; \texttt{RandomBrightnessContrast(p=0.5)}
    для адаптації до варіацій освітлення;
    \texttt{HueSaturationValue(hue\_shift\_limit=15, p=0.3)} для робастності до
    кольорового балансу.
    \item \textbf{Запуск тренування.} Виклик
    \texttt{model.train(data, epochs=100, imgsz=640, batch=16, patience=20,
    optimizer="AdamW", lr0=1e-3, weight\_decay=5e-4)} з валідацією раз на епоху.
    Параметр \texttt{patience=20} активує ранню зупинку при відсутності покращення
    метрики mAP50 протягом 20 епох.
    \item \textbf{Експорт результуючих ваг.} Найкращі ваги (за метрикою mAP50-95)
    зберігаються у файл \texttt{ball\_best\_detec\_server\_v2.pt} та переносяться
    до каталогу \texttt{models/}.
\end{enumerate}

Конфігурація набору даних задається через файл
\texttt{server\_data.yaml}, що містить шляхи до тренувальної,
валідаційної та тестової частин, перелік класів (\texttt{names: [ball]}) та
кількість класів (\texttt{nc: 1}). Сам набір даних обсягом близько 18~000 анотованих
кадрів отриманий шляхом ручної розмітки 5 окремих волейбольних записів та доповнений
синтетичними даними з відкритого датасету Volleyball Detection 5.

\subsection{Конфігураційні параметри та залежності}

Усі обчислювально-критичні константи системи виведені у окремий перелік для
зручності тонкого налаштування під специфічні умови зйомки. Зведений перелік
наведено у таблиці~\ref{tbl:config_params}.

\begin{table}[!htbp]
\fontsize{11pt}{13pt}\selectfont
\caption{Зведений перелік ключових конфігураційних параметрів системи}
\label{tbl:config_params}
\centering
\begin{tblr}{
    width = \textwidth,
    colspec = {| X[2,l,m] | X[1.2,c,m] | X[1,c,m] | X[3,l,m] |},
    hlines,
    vlines,
    row{1} = {font=\bfseries, bg=lightgray!20},
}
    Параметр & Поточне значення & Одиниці & Призначення \\
    \texttt{dt} & 0.02 & с & Крок дискретизації (1/FPS) \\
    \texttt{max\_age} & 50 & кадри & Поріг видалення неактивного треку \\
    \texttt{min\_hits} & 3 & кадри & Поріг підтвердження треку \\
    \texttt{gating\_threshold} & 11.34 & --- & $\chi^2$-поріг (df=3, P=99\%) \\
    \texttt{ball\_diameter} & 0.21 & м & Стандартний діаметр м'яча FIVB \\
    $\alpha$ & 0.1 & --- & Розкид сигма-точок \\
    $\beta$ & 2.0 & --- & Параметр гаусового розподілу \\
    $\kappa$ & 1.0 & --- & Допоміжний параметр сигма-точок \\
    $q_{ballistic}$ & 0.1 & м$^2$/с$^4$ & Дисперсія Q для балістики \\
    $q_{hit}$ & 100.0 & м$^2$/с$^4$ & Дисперсія Q для удару \\
    $q_{bounce}$ & 5.0 & м$^2$/с$^4$ & Дисперсія Q для відскоку \\
    $\sigma_x^2 = \sigma_y^2$ & 0.01 & м$^2$ & Дисперсія R по горизонталі \\
    $\sigma_z^2$ & 0.04 & м$^2$ & Дисперсія R по глибині \\
    $\varepsilon$ & 0.75 & --- & Коефіцієнт реституції відскоку \\
    $\mu_{friction}$ & 0.85 & --- & Коефіцієнт тангенційного тертя \\
    $k_{drag}$ & 0.0314 & 1/(с$\cdot$м) & Зведений коефіцієнт аеродинаміки \\
\end{tblr}
\end{table}

\textbf{Програмні залежності.} Проект розгортається у контрольованому середовищі
Anaconda на основі manifesta \texttt{packages.txt}. Ключові бібліотеки із вказаними
версіями:
\begin{itemize}
    \item \texttt{numpy} 1.26.4 --- базові матричні операції;
    \item \texttt{scipy} 1.13.1 --- алгоритми наукових обчислень (Cholesky, Hungarian,
    multivariate normal);
    \item \texttt{opencv-python} 4.10.0 --- комп'ютерний зір (solvePnP, Rodrigues);
    \item \texttt{ultralytics} 8.3.40 --- API для моделей \gls{YOLO};
    \item \texttt{albumentations} 1.4.18 --- аугментації навчальної вибірки;
    \item \texttt{torch} 2.4.1 + \texttt{torchvision} 0.19.1 --- фреймворк глибокого
    навчання (бекенд для Ultralytics);
    \item \texttt{matplotlib} 3.9.2 --- візуалізація траєкторій та метрик;
    \item \texttt{jupyter} 1.0.0 --- інтерактивне середовище для експериментальних
    ноутбуків.
\end{itemize}

\textbf{Вимоги до апаратного забезпечення.} Мінімальна конфігурація для роботи у
режимі реального часу: процесор з 6 фізичними ядрами (Intel i5-12500 або AMD Ryzen
5 5600 і вище), 16 ГБ оперативної пам'яті, відеокарта з 8 ГБ VRAM та підтримкою
CUDA 12.0 (NVIDIA RTX 3060 і вище). Рекомендована конфігурація: відеокарта рівня
NVIDIA RTX 4070 або вище, 32 ГБ оперативної пам'яті, NVMe SSD для зберігання
вхідних відеопотоків.

\subsection{Інструкція з користування програмним комплексом}

Запуск повного циклу обробки відеозапису волейбольного матчу здійснюється у
декілька послідовних кроків з командного рядка:

\begin{enumerate}
    \item Активація віртуального середовища: \texttt{conda activate volleyball\_analysis}.
    \item Перевірка коректності тренованих ваг моделі: повинен існувати файл
    \texttt{training\_models/models/ball\_best\_detec\_server\_v2.pt}.
    \item Запуск інтерактивної версії конвеєра через Jupyter Notebook:
    \texttt{jupyter notebook ball\_3d\_analysis/IMM\_UKF\_modeling.ipynb}.
    \item Послідовне виконання комірок ноутбука: завантаження моделі та відео,
    калібровка камери, попередня детекція з збереженням у JSON, ініціалізація
    трекера, поетапна обробка кадрів, побудова візуалізації траєкторій.
    \item Аналіз отриманих результатів: розрахунок метрик гладкості траєкторій,
    розподілу режимів IMM, кількості перемикань ідентифікаторів треків.
\end{enumerate}

Альтернативний пакетний режим обробки, призначений для масових експериментальних
прогонів на серіях коротких відеосегментів з автоматичним розрахунком метрик
якості, є частиною подальшого розвитку системи і буде описаний у наступних версіях
документації.

\begin{unnumbered}

\section*{Висновки до розділу 3}

У межах третього розділу здійснено повну програмну реалізацію розробленого
тривимірного аналітичного конвеєра у вигляді чотирьох логічно та структурно
розмежованих модулів на мові Python. Кожен модуль інкапсулює одну фундаментальну
функцію системи: тривимірне геометричне позиціонування, нелінійну стохастичну
фільтрацію, асоціацію даних з керуванням життєвим циклом треків та навчання
згорткової нейромережі детекції. Описаний рівень модульності забезпечує можливість
ізольованого модифікування та тестування окремих компонентів без впливу на стабільність
сусідніх підсистем, що є необхідною передумовою подальшого масштабування комплексу.

Алгоритмічне ядро системи реалізоване через три основні класи (\texttt{Unscented\-Kalman\-Filter},
\texttt{IMMEstimator}, \texttt{MerweScaledSigmaPoints}) та два класи високого рівня
(\texttt{Track}, \texttt{IMMTracker}), що разом утворюють завершену реалізацію
гібридної архітектури взаємодіючих множинних моделей з адаптивною матрицею переходу
режимів. Особлива увага була приділена коректній підтримці випадків відсутності
вимірювань (метод \texttt{mark\_missed}) та захисту проти числових аномалій
(виродженої матриці інновації, експоненціально великих швидкостей через шум
детекції), що гарантує безперебійну роботу фільтра навіть на складних відеосегментах
з тривалими окклюзіями м'яча.

Геометричний модуль на основі алгоритму \gls{PnP} та принципу зворотного проектування
демонструє інженерну простоту та прозорість: уся каліброчна процедура зведена до
двох публічних функцій з чітко окресленими вхідними та вихідними інтерфейсами. Такий
дизайн дозволяє повторно використовувати модуль для будь-якого нового запису після
ручного визначення піксельних координат чотирьох реперних точок розмітки поля.

Програмний комплекс модуля навчання моделі \gls{YOLO}26 використовує сучасний
фреймворк Ultralytics з широким набором візуальних аугментацій, що моделюють
типові спотворення волейбольного відеопотоку (рухове розмиття, шум сенсора,
варіації освітлення, кольорового балансу). Завдяки цьому навчена модель демонструє
підвищену стійкість до екстремальних умов реальної гри без необхідності збирати
непропорційно великий ручний набір даних.

Сукупність розроблених модулів становить функціонально повний інженерний інструментарій,
готовий до експериментальної верифікації та подальшого розширення додатковими
підсистемами розрізання відеосегментів, автоматичного обчислення метрик якості та
формування звітів. Програмний комплекс задовольняє вимогам реального часу на типовому
користувацькому обладнанні і повністю відповідає теоретичним рішенням, обґрунтованим у
попередніх розділах роботи.

\end{unnumbered}
